\chapter{調査計画：記録すべき変数と事前登録}
\label{ch:next}
%============================================================

\section{調査の優先順位}

対象を族2 と族5 に絞る。選定は演繹で導ける。

\paragraph{$\phi_{\mathrm{prod}}$ は $\Phi$ に依存しない。}

式\eqref{eq:surplus}の三分解のうち $\phi_{\mathrm{prod}}=(p^{\ast}-c(\delta))\delta$ は
限界費用と競争価格で決まり、履行から決済への写像の形を含まない。
同一の $f$ と同一の限界費用に対して $\Phi$ だけを差し替えても、
変わるのは $\kappa$ の符号と必要な運転資本であって生産余剰ではない
（表\ref{tab:same-f-different-phi}）。

\textbf{$\Phi$ の形が余剰に効くのは残る二項においてである。}
枠組みを検証するには、この二項をそれぞれ単独に含む $\Phi$ を選べばよい。

\paragraph{二項は別の自由度から生じる。}

$\phi_{\mathrm{cog}}=p\cdot\E[\dbar-\delta]$ が正となりうるのは、
契約が上限 $\dbar$ を明示する場合、すなわち $\pi$ が $\delta$ を参照しない場合に限られる
（注\ref{prop:breakage}）。これは自由度(5)である。
加えて、乖離を取引頻度 $\nu$ に結びつけるには自由度(6)の反復性を要する。
単発では $\dbar-\delta$ が一回の実現でしかなく分布を持たないため、
$\nu$ との関係を問えない。
\textbf{自由度(5)と(6)がともに立つ族は族2}であり、
第\ref{part:catalog}部はこれを認知余剰の主要な所在として定めている。

$\phi_{\mathrm{barg}}=(p-p^{\ast})\delta$ はゼロサムの移転であり、
その配分は $\pi=h(y)$ の $y$ が検証可能かどうかで決まる。これは自由度(4)である。
表\ref{tab:family-assignment-order}の判定順序において、
\textbf{自由度(4)で判定される族は族5のみ}である。

\paragraph{二族は象限図の対角にある。}

自由度(1)と(2)が支配的であり、この二軸が図\ref{fig:quadrant}の四分割を与える。
族2 は前受・定額型（$\pi$ が $\delta$ と独立、乖離しうる）、
族5 は後払・連動型（$\pi$ が $\delta$ に連動、乖離しない）であって、
\textbf{二軸の値がいずれも逆になる}。
どちらか一族だけを選べば、動かせるのは二軸のうち片方だけである。

\paragraph{何を検証できるか。}

族2 は認知余剰（式\eqref{eq:cog}）と容量制約の命題を直接検証できる。
特に、同一の定額会員制でありながら容量制約の有無で二分される点は、
外形的に似た $\Phi$ が異なる耐久性を持つことの最も明快な事例になる。

族5 は可測性条件（第\ref{sec:info}章）を直接検証できる。
成果報酬とコストプラスが逆向きのリスク移転であるという主張は、
契約形態と観測可能性の対応として検証可能である。

\paragraph{並べる理由。}

二族は制度的な根拠が異なるため、データの得られ方も異なる。
族2 は資金決済法の登録制により全数フレームが存在し、
族5 は職業安定法の情報提供義務により個社の届出が公開されている。
\textbf{前者は網羅的な枠を与え、後者は個体の変数を与える。}

一族だけを見ると、検証が進まなかったときに、
その原因を\textbf{フレームの不備に帰すべきか変数の不備に帰すべきか決まらない}。
根拠の異なる二族を並べれば、どちらの側で詰まったのかを切り分けられる。
二族を対象とする理由はここにもある。

\section{記録すべき変数}

各 $\Phi$ について、最低限次を記録する。

\begin{center}
\begin{tabularx}{\textwidth}{@{}lX@{}}
\toprule
変数 & 内容 \\
\midrule
$\operatorname{sgn}\kappa$ & 決済が履行に先行するか後行するか \\
$\CCC$ & 定常状態が成立する場合のみ。成立しない場合はその旨を記録 \\
依存形式 & $c$ / $p\delta$ / $F+p\delta$ / $h(y)$ のいずれか \\
$\dbar$ と $\delta$ の分離 & 分離の有無、および乖離の推定値 \\
容量制約 & 有無、および制約が拘束的かどうか \\
$\nu$ & 取引頻度（式\eqref{eq:decay}） \\
可測性 & 努力・結果のいずれが検証可能か \\
支払主体 & 受益者本人か、第三者か \\
$\phi$ の推定内訳 & $\phi_{\mathrm{prod}}$ / $\phi_{\mathrm{barg}}$ / $\phi_{\mathrm{cog}}$ の定性的配分 \\
\bottomrule
\end{tabularx}
\captionof{table}{$\Phi$ ごとに最低限記録すべき変数の一覧。}
\label{tab:variables-to-record}
\end{center}

\section{事前登録すべき事項}

第\ref{part:method}部の結論に従い、調査開始前に以下を確定し記録する。

\begin{enumerate}[label=(\arabic*)]
\item 母集団の定義（基準時点、地理的範囲、業種範囲、規模下限）
\item 定常状態とみなす閾値（$r$ の安定性の判定基準）
\item 失敗事例の収集方法と、収集不能な範囲
\item APC のうちゼロと仮定する効果、およびその根拠
\item 一次資料と回顧的記述の区別の基準
\end{enumerate}

これらを事後に決めると、第\ref{sec:survivorship}節および第\ref{sec:disclosure}節の
選択が結論に混入する。

